% Arquivo principal. Compilador: pdfLaTeX. Nao requer BibTeX ou Biber.
\documentclass[12pt,a4paper]{article}

% true: exibe orientacoes em portugues e todos os titulos de trabalho.
% false: oculta orientacoes e espacos vazios das partes ainda nao escritas.
\newif\iforientacoes
\orientacoestrue
% Permite verificar o modo limpo pela linha de comando, sem editar o arquivo.
\ifdefined\SemOrientacoes\orientacoesfalse\fi

% Opcoes editoriais de trabalho, nao medidas impostas pela revista:
% A4, 12 pt, margens de 2,5 cm, entrelinha 1,5, Latin Modern.
\usepackage[T1]{fontenc}
\usepackage[utf8]{inputenc}
\usepackage{lmodern}
\IfFileExists{portuguese.ldf}{%
  \usepackage[portuguese,english]{babel}
}{%
  \usepackage[english]{babel}
  % Se os padroes portugueses nao estiverem instalados, nao aplicar
  % hifenizacao inglesa aos trechos em portugues.
  \babelprovide[import,hyphenrules=nohyphenation]{portuguese}
}
\usepackage[margin=2.5cm]{geometry}
\usepackage{microtype}
\usepackage{setspace}
\onehalfspacing
\usepackage[autostyle=false,style=english]{csquotes}
\usepackage{enumitem}
\usepackage{xcolor}
\usepackage{etoolbox}
\usepackage{xparse}
\usepackage{needspace}
\usepackage[most]{tcolorbox}
\usepackage{titlesec}
\usepackage{xurl}
\usepackage[hidelinks,unicode]{hyperref}
\urlstyle{same}
\hypersetup{
  pdftitle={How to talk about the Anthropocene? A Science-bar essay towards ancestral futures},
  pdfauthor={},
  pdfsubject={Essay manuscript},
  pdfkeywords={Anthropocene, Bruno Latour, common world, ancestral futures}
}

\setlength{\parindent}{1.25cm}
\setlength{\parskip}{0pt}
\setlength{\emergencystretch}{2em}
\clubpenalty=10000
\widowpenalty=10000
\displaywidowpenalty=10000
\raggedbottom
\titleformat{\section}{\large\bfseries}{\thesection.}{0.6em}{}
\titlespacing*{\section}{0pt}{2em}{0.8em}
\renewcommand{\thesection}{\Roman{section}}
\setlist[enumerate]{leftmargin=1.7em,itemsep=0.25em,topsep=0.4em}
\setlist[itemize]{leftmargin=1.4em,itemsep=0.25em,topsep=0.4em}

% Conteudo editorial: completamente removido no modo sem orientacoes.
\definecolor{guideline}{HTML}{BDC6CB}
\definecolor{guidefill}{HTML}{F4F6F7}
\NewDocumentEnvironment{orientacao}{m +b}{%
  \iforientacoes
    \begin{tcolorbox}[enhanced,breakable,colback=guidefill,colframe=guideline,
      boxrule=0.35pt,arc=0pt,left=9pt,right=9pt,top=7pt,bottom=7pt,
      before skip=1em,after skip=1em,
      title={#1},colbacktitle=guidefill,coltitle=black,
      fonttitle=\normalfont\sffamily\bfseries\small,
      fontupper=\normalfont\small]
      \selectlanguage{portuguese}\raggedright
      \setstretch{1.13}\setlength{\parindent}{0pt}\setlength{\parskip}{0.4em}
      #2
    \end{tcolorbox}
  \fi
}{}

% Atos: narrativa ficcional, sem aspas de citacao de fonte externa.
\newcounter{ato}[section]
\NewDocumentEnvironment{ato}{m}{%
  \refstepcounter{ato}
  \par\addvspace{1.2em}
  \noindent\begin{minipage}{\linewidth}
  \noindent\textbf{Act \Roman{ato}. #1}\par\smallskip
  \begingroup\itshape\setstretch{1.2}
  \setlength{\parindent}{0pt}\setlength{\parskip}{0.35em}
}{\par\endgroup\end{minipage}\par\medskip}

% Uma chave centraliza a nota curta e a referencia completa.
% O argumento opcional de \fonte recebe o localizador ja formatado.
\newcommand{\declararobra}[3]{%
  \csdef{curta@#1}{#2}\csdef{completa@#1}{#3}}
\newcommand{\usarobra}[1]{\csgdef{usada@#1}{1}}
\newcommand{\refobra}[1]{%
  \ifcsdef{curta@#1}{\usarobra{#1}\csuse{curta@#1}}{%
    \PackageError{sciencebar}{Referencia desconhecida: #1}{Cadastre a chave em referencias.tex.}}}
\NewDocumentCommand{\fonte}{O{} m}{%
  \footnote{\refobra{#2}\ifstrempty{#1}{}{, #1}.}}
\newcommand{\entradaobra}[1]{%
  \iforientacoes
    \par\noindent\hangindent=1.25cm\hangafter=1\csuse{completa@#1}\par\medskip
  \else
    \ifcsdef{usada@#1}{%
      \par\noindent\hangindent=1.25cm\hangafter=1\csuse{completa@#1}\par\medskip}{}
  \fi}

% As aspas sao visiveis e deliberadas; o texto explica sua funcao.
\newcommand{\anthropocene}{\enquote{Anthropocene}}
\newcommand{\sciencebar}{Science-bar}

% Espaco para atos ainda nao escritos nas partes II e III.
% Ao escrever, substituir \atoplanejado por \begin{ato} ... \end{ato}.
\newcommand{\atoplanejado}[2]{%
  \begin{orientacao}{Ato #1 --- a escrever}#2\end{orientacao}}

% Referencias manuais no modelo MLA 9 com notas adotado por Das Questoes.
% Nao requer BibTeX/Biber. Notas curtas sem ano; localizador opcional.
% Para uma nova obra: declarar a chave e inclui-la, em ordem alfabetica,
% na lista de \imprimirreferencias ao final deste arquivo.

\declararobra{crutzen2002}
  {Crutzen, \enquote{Geology of Mankind}}
  {Crutzen, Paul J. \enquote{Geology of Mankind.} \emph{Nature},
   vol.~415, 2002, p.~23. \url{https://doi.org/10.1038/415023a}.}

\declararobra{crutzen-stoermer2000}
  {Crutzen and Stoermer, \enquote{The \enquote{Anthropocene}}}
  {Crutzen, Paul J., and Eugene F. Stoermer. \enquote{The \enquote{Anthropocene}.}
   \emph{Global Change Newsletter}, no.~41, May 2000, pp.~17--18.
   \emph{Max Planck Institute for Chemistry},
   \url{https://www.mpic.de/3864697/the-anthropocene}.}

\declararobra{cunha2017}
  {Cunha, \emph{Cultura com Aspas}}
  {Cunha, Manuela Carneiro da. \emph{Cultura com Aspas e Outros Ensaios}.
   Ubu Editora, 2017.}

\declararobra{foucault2013}
  {Foucault, \emph{O Corpo Utópico, As Heterotopias}}
  {Foucault, Michel. \emph{O Corpo Utópico, As Heterotopias}.
   Translated by Salma Tannus Muchail, n-1 edições, 2013.}

\declararobra{haraway2015}
  {Haraway, \enquote{Anthropocene, Capitalocene, Plantationocene, Chthulucene}}
  {Haraway, Donna. \enquote{Anthropocene, Capitalocene, Plantationocene,
   Chthulucene: Making Kin.} \emph{Environmental Humanities}, vol.~6,
   2015, pp.~159--65. \url{https://doi.org/10.1215/22011919-3615934}.}

\declararobra{haraway2016}
  {Haraway, \enquote{Tentacular Thinking}}
  {Haraway, Donna. \enquote{Tentacular Thinking: Anthropocene,
   Capitalocene, Chthulucene.} \emph{e-flux Journal}, no.~75, Sept.~2016,
   \url{https://www.e-flux.com/journal/75/67125/tentacular-thinking-anthropocene-capitalocene-chthulucene/}.}

\declararobra{iugs2024}
  {IUGS and ICS, \enquote{The Anthropocene}}
  {International Union of Geological Sciences, and International Commission
   on Stratigraphy. \enquote{The Anthropocene: IUGS-ICS Statement.}
   20 Mar.~2024,
   \url{https://www.iugs.org/wp-content/uploads/2024/03/Anthropocene_extended_IUGS-ICS_Statement-1.pdf}.}

\declararobra{krenak2022}
  {Krenak, \emph{Futuro Ancestral}}
  {Krenak, Ailton. \emph{Futuro Ancestral}. Companhia das Letras, 2022.}

\declararobra{latour2004}
  {Latour, \enquote{How to Talk About the Body?}}
  {Latour, Bruno. \enquote{How to Talk About the Body? The Normative
   Dimension of Science Studies.} \emph{Body \& Society}, vol.~10,
   nos.~2--3, 2004, pp.~205--29.
   \url{https://doi.org/10.1177/1357034X04042943}.}

\declararobra{latour1994}
  {Latour, \emph{Jamais Fomos Modernos}}
  {Latour, Bruno. \emph{Jamais Fomos Modernos: Ensaio de Antropologia
   Simétrica}. Translated by Carlos Irineu da Costa, Editora 34, 1994.}

\declararobra{latour2020}
  {Latour, \emph{Onde Aterrar?}}
  {Latour, Bruno. \emph{Onde Aterrar? Como Se Orientar Politicamente
   no Antropoceno}. Translated by Marcela Vieira, Bazar do Tempo, 2020.}

\declararobra{lewis-maslin2015}
  {Lewis and Maslin, \enquote{Defining the Anthropocene}}
  {Lewis, Simon L., and Mark A. Maslin. \enquote{Defining the Anthropocene.}
   \emph{Nature}, vol.~519, 2015, pp.~171--80.
   \url{https://doi.org/10.1038/nature14258}.}

\declararobra{malm-hornborg2014}
  {Malm and Hornborg, \enquote{The Geology of Mankind?}}
  {Malm, Andreas, and Alf Hornborg. \enquote{The Geology of Mankind?
   A Critique of the Anthropocene Narrative.} \emph{The Anthropocene
   Review}, vol.~1, no.~1, 2014, pp.~62--69.
   \url{https://doi.org/10.1177/2053019613516291}.}

\declararobra{martins2021}
  {Martins, \emph{Performances do Tempo Espiralar}}
  {Martins, Leda Maria. \emph{Performances do Tempo Espiralar:
   Poéticas do Corpo-Tela}. Cobogó, 2021.}

\declararobra{mauss2017}
  {Mauss, \emph{Sociologia e Antropologia}}
  {Mauss, Marcel. \emph{Sociologia e Antropologia}.
   Translated by Paulo Neves, Ubu Editora, 2017.}

\declararobra{ruddiman2003}
  {Ruddiman, \enquote{The Anthropogenic Greenhouse Era}}
  {Ruddiman, William F. \enquote{The Anthropogenic Greenhouse Era
   Began Thousands of Years Ago.} \emph{Climatic Change}, vol.~61,
   2003, pp.~261--93. \url{https://doi.org/10.1023/B:CLIM.0000004577.17928.fa}.}

\declararobra{steffen2015}
  {Steffen et al., \enquote{The Trajectory of the Anthropocene}}
  {Steffen, Will, et al. \enquote{The Trajectory of the Anthropocene:
   The Great Acceleration.} \emph{The Anthropocene Review}, vol.~2,
   no.~1, 2015, pp.~81--98. \url{https://doi.org/10.1177/2053019614564785}.}

\declararobra{awg}
  {Subcommission on Quaternary Stratigraphy, \enquote{Working Group}}
  {Subcommission on Quaternary Stratigraphy. \enquote{Working Group on
   the \enquote{Anthropocene}.} \emph{International Commission on
   Stratigraphy}, \url{https://quaternary.stratigraphy.org/working-groups/anthropocene}.
   Accessed 5 Oct.~2026.}

\newcommand{\imprimirreferencias}{%
  \begingroup\setstretch{1.15}\setlength{\parindent}{0pt}
  \entradaobra{crutzen2002}
  \entradaobra{crutzen-stoermer2000}
  \entradaobra{cunha2017}
  \entradaobra{foucault2013}
  \entradaobra{haraway2015}
  \entradaobra{haraway2016}
  \entradaobra{iugs2024}
  \entradaobra{krenak2022}
  \entradaobra{latour2004}
  \entradaobra{latour1994}
  \entradaobra{latour2020}
  \entradaobra{lewis-maslin2015}
  \entradaobra{malm-hornborg2014}
  \entradaobra{martins2021}
  \entradaobra{mauss2017}
  \entradaobra{ruddiman2003}
  \entradaobra{steffen2015}
  \entradaobra{awg}
  \endgroup}


\title{How to talk about the Anthropocene?\\[0.35em]
  \large A Science-bar essay towards ancestral futures}
\author{} % Avaliacao anonima: preencher dados_autoria.tex separadamente.
\date{}

\begin{document}
\maketitle
\vspace{-2.5em}

\begin{orientacao}{Como trabalhar neste arquivo}
Este documento retoma a proposta de sete atos para \textit{Crisis}.
Os atos estão redigidos em inglês; os parágrafos acadêmicos são inícios
provisórios de redação, a desenvolver. As caixas em português contêm
orientações editoriais. A hipótese do ensaio aparece como proposta de quem
escreve, distinguida das ideias atribuídas às fontes.

Altere \texttt{\string\orientacoestrue} para
\texttt{\string\orientacoesfalse} em \texttt{main.tex} para ocultar as
caixas e os espaços ainda vazios. Essa operação produz uma leitura sem
anotações; não transforma o rascunho em artigo concluído. Os resumos
também são versões provisórias. Veja \texttt{README.md} para compilar,
inserir notas e preparar o arquivo de submissão.
\end{orientacao}

% Reescrever ambos os resumos ao terminar o ensaio; manter equivalencia.
% O tamanho e a quantidade de palavras-chave abaixo sao escolhas editaveis.
\begin{orientacao}{Resumos provisórios}
As duas versões abaixo enunciam o projeto, sem apresentar resultados como
concluídos. Ajuste-as depois de escrever as três partes, sobretudo a
definição de Cecinepá. Para manter a avaliação anônima, examine também
referências autobiográficas, agradecimentos e a descrição desse espaço.
\end{orientacao}

\Needspace{5\baselineskip}
\noindent\textbf{Abstract}\par
\noindent
This essay approaches contemporary systemic crises through Anthropocene
talk: the practices through which a changing world is named, described,
and made available for collective action. A fictional sciences bar runs
alongside the argument, following a Purifier whose efforts to conceal
mediations progressively alter the conditions of speech. Drawing on Bruno
Latour and the controversies surrounding the Anthropocene, the essay
proposes quotation marks as a sign of situated borrowing and translation.
It develops this proposal through Science-bar and Cecinepá, bringing
Latour's common world into conversation with Manuela Carneiro da Cunha's
account of culture in quotation marks, Michel Foucault's heterotopias,
Marcel Mauss's total social fact, Ailton Krenak's ancestral futures, and
Leda Maria Martins's spiral time. The guiding question concerns how to
articulate differences and dependencies while composing conditions for
living together.

\medskip
\noindent\textbf{Keywords:} Anthropocene; quotation marks; Science-bar;
ancestral futures; common world.

\begin{otherlanguage}{portuguese}
\Needspace{5\baselineskip}
\bigskip
\noindent\textbf{Resumo}\par
\noindent
Este ensaio aborda as crises sistêmicas contemporâneas a partir do
\textit{Anthropocene talk}: as práticas pelas quais um mundo em
transformação é nomeado, descrito e tornado disponível à ação coletiva.
Um bar de ciências ficcional acompanha o argumento, seguindo um
Purificador cujos esforços para ocultar mediações alteram progressivamente
as condições da fala. Em diálogo com Bruno Latour e com as controvérsias
em torno do Antropoceno, o ensaio propõe as aspas como sinal de empréstimo
situado e tradução. Desenvolve essa proposta por meio de Science-bar e
Cecinepá, aproximando o mundo comum de Latour da cultura com aspas de
Manuela Carneiro da Cunha, das heterotopias de Michel Foucault, do fato
social total de Marcel Mauss, dos futuros ancestrais de Ailton Krenak e do
tempo espiralar de Leda Maria Martins. A pergunta orientadora concerne a
como articular diferenças e dependências enquanto se compõem condições
de vida em comum.

\medskip
\noindent\textbf{Palavras-chave:} Antropoceno; aspas; Science-bar;
futuros ancestrais; mundo comum.
\end{otherlanguage}

\section{Crisis}
\label{sec:crisis}

\begin{orientacao}{Percurso da primeira parte}
Perguntar $\rightarrow$ seguir o nome $\rightarrow$ disputar seu começo
$\rightarrow$ distribuir responsabilidades $\rightarrow$ reconhecer as
regras de purificação $\rightarrow$ experimentar suas consequências na
fala $\rightarrow$ situar o empréstimo por meio das aspas.

O primeiro ato conserva a formulação fornecida por você. Os seis seguintes
retomam a primeira proposta. Cada bloco acadêmico começa com um rascunho
em inglês, seguido de tarefas para ampliá-lo. Evite transformar a estória
em ilustração com equivalências fixas: o tapete, a poeira, os corpos e os
documentos podem produzir conexões novas ao longo do ensaio.
\end{orientacao}

\begin{ato}{The question at the counter}
The writer walks into a bar; it’s a sciences bar, that’s where scientists,
disciplines and epistemologies hang out and drink together; the writer
then hits the counter and screams his lungs out for a fine answer to his
fair question: how the hell can one talk about the Anthropocene?!
\end{ato}

% Inicio academico provisorio. Desenvolver com uma situacao efetiva de escrita.
The question reaches the counter before the writer has learned how to
listen in this room. It already presupposes an answer that could settle
the matter, as though speaking more forcefully would make the world more
available. Taking up Latour's approach to body talk, I ask what becomes
perceptible through the ways we speak of the Anthropocene, and what those
ways of speaking leave us unable to register.\fonte[pp.~205--08]{latour2004}
This requires situating the activity of writing and reading among the
dependencies it seeks to describe. The place from which the question is
asked is part of the question.

\begin{orientacao}{Pós-ato I --- situar a escrita e abrir o problema}
\textbf{Função.} Localizar a escrita-leitura em uma situação concreta e
apresentar o deslocamento para o \textit{Anthropocene talk}.
\begin{enumerate}
\item Abra com uma experiência real do lugar e do período em que você
escreve: ar, água, calor, trabalho, infraestrutura ou outra dependência
efetivamente observada. Descreva a conexão que permite passar dessa
situação às crises sistêmicas. Não invente uma vivência para preencher
o modelo; acrescente uma fonte quando fizer uma afirmação empírica.
\item Explique o deslocamento inspirado no \textit{body talk}: interessa
seguir como o Antropoceno é mobilizado em relatos, instrumentos,
instituições e controvérsias. A pergunta inclui quem fala, de onde, com
quais condições de escuta e com quais efeitos de conhecimento.
\item Explore a confiança do escritor em obter uma resposta definitiva.
O gesto de bater no balcão e a referência aos pulmões antecipam sua
implicação material no bar. Ele também terá de aprender a ser afetado;
não é um narrador que observa tudo de fora.
\end{enumerate}
\textbf{Passagem.} Uma palavra chega como resposta à pergunta, mas é
preciso acompanhar sua procedência e seus modos de circulação.
\end{orientacao}

\begin{ato}{The travelling label}
The bartender reaches for a bottle labelled ANTHROPOCENE.

“This one has travelled,” she says. “Laboratories, newspapers, parliaments.
People keep ordering it before agreeing on the recipe.”

The writer turns the bottle in his hands. Beneath its label, the edges of
older labels show through.

A figure wearing a badge that reads PURIFIER brushes a trail of grit
beneath the carpet.

“Just the packaging,” he says.

The bartender pours. The writer watches something settle at the bottom
of his glass.
\end{ato}

The name arrives with a history of observations, instruments, and claims
to authority. In 2000, Paul Crutzen and Eugene Stoermer proposed
Anthropocene for a present increasingly transformed by human activities;
their short intervention also assembled earlier accounts of humanity as
a geological force.\fonte{crutzen-stoermer2000}
Crutzen's 2002 article extended the proposal in \textit{Nature}.\fonte{crutzen2002}
Following this trajectory involves more than recovering a word's first
appearance. As the name moves between practices, it gathers different
objects of concern and different expectations about what its use can
accomplish. A circulating label does not carry an unchanged world inside
it. Its translations become part of what this essay must learn to follow.

\begin{orientacao}{Pós-ato II --- fazer uma genealogia da circulação}
\textbf{Função.} Apresentar uma genealogia seletiva que explique a força
do termo e as diferenças entre seus usos.
\begin{enumerate}
\item Recupere os antecedentes reunidos por Crutzen e Stoermer:
Marsh (1864), Stoppani (1873, \textit{Anthropozoic}) e Vernadsky
(1926, biosfera). Eles não são três invenções do mesmo termo. Se você
não consultar essas obras diretamente, indique que a reconstrução é
feita por meio do artigo de 2000.
\item Situe o uso informal por Stoermer nos anos 1980, conforme a
genealogia de Haraway, e diferencie-o da intervenção conjunta de 2000
e da projeção alcançada em 2002.
\item Acompanhe a passagem entre ciência do sistema Terra, estratigrafia,
humanidades e debate público. Escolha exemplos documentados; mostre
o que muda em cada tradução, em vez de apresentar uma lista de nomes.
\end{enumerate}
\textbf{Apoio.} A genealogia de Haraway orienta a seleção de interlocutores
e a pergunta sobre as histórias que cada nome permite contar.
Veja \refobra{haraway2016}.

\textbf{Simetria.} Os rótulos sobrepostos sugerem histórias diferentes;
o sedimento no copo encaminha a investigação sobre a marca material de
um começo. A primeira varrida já reduz uma mediação a mero invólucro.
\end{orientacao}

\begin{ato}{A beginning under examination}
A geologist places a narrow cylinder of sediment beside the drink.
Someone points to a layer; someone else asks how far its trace can be
followed.

Dates spread across the counter.

“When does it begin?” asks the writer.

“What, exactly?” replies the geologist.

A notice arrives: THE PROPOSED EPOCH HAS NOT BEEN APPROVED.

The Purifier folds it into a wedge and steadies the counter. For a moment,
the glasses stop rattling.

A draught lifts another thread of dust from the carpet.
\end{ato}

The question of a beginning changes with the object being dated. An
account of human transformations of environments, a diagnosis of change
in the Earth system, and a proposal for a formal unit of geological time
do not necessarily require the same boundary. The Anthropocene Working
Group proposed a boundary in 1952, using the sedimentary record at
Crawford Lake.\fonte{awg}
In March 2024, the IUGS and ICS endorsed the rejection of that formal
epoch proposal.\fonte{iugs2024}
The decision concerns geological classification; it does not establish
that the transformations gathered under the name are unreal. The
Purifier's use of the notice as a stabilising wedge performs a different
operation: it turns a decision about classification into a promise that
the room can remain as it is.

\begin{orientacao}{Pós-ato III --- distinguir os objetos da controvérsia}
\textbf{Função.} Relacionar propostas de início a critérios e explicações
diferentes. Em estratigrafia, a proposta era de uma \textit{época}, não de
uma nova \textit{era}. Evite intercambiar essas categorias técnicas.
\begin{enumerate}
\item Compare, em prosa, quatro conjuntos de propostas: agricultura
antiga e alterações de gases de efeito estufa em Ruddiman; fim do
século XVIII em Crutzen e Stoermer; o marco de 1610 discutido por
Lewis e Maslin; e a transformação de meados do século XX associada
à Grande Aceleração. Para 1610, explicite a hipótese que relaciona
colonização, mortalidade indígena, regeneração florestal e queda de
CO\textsubscript{2}; apresente-a como uma proposta discutida.
\item Diferencie a Grande Aceleração do marcador estratigráfico
específico. O plutônio no registro de Crawford Lake serve a uma operação
de correlação e delimitação; não é uma explicação suficiente de todas
as crises. A proposta de 2023 do grupo de trabalho, criado em 2009,
chegou ao procedimento de avaliação rejeitado em 2024.
\item Explique o alcance institucional da decisão e preserve a
legitimidade das questões estratigráficas. O geólogo da estória não
precisa ser confundido com o Purificador. O gesto deste é converter
um resultado localizado em imunidade para todo o bar.
\end{enumerate}
\textbf{Fontes para o desenvolvimento.}
\refobra{ruddiman2003}; \refobra{lewis-maslin2015};
\refobra{steffen2015}. Antes da submissão, confira novamente o estado
da classificação na ICS: a controvérsia e a pesquisa continuaram após
a decisão de 2024. Uma data de formalização não substitui a história
das responsabilidades que o próximo ato interroga.
\end{orientacao}

\begin{ato}{The bill addressed to humanity}
The bill arrives addressed to HUMANITY.

“Shall we split it equally?” asks a guest ordering another round.

“I have been carrying the glasses,” says a woman by the door.

Someone opens the ledger. Its pages record who ordered, who served, and
who owns the bar.

At one table, CAPITALOCENE is written across a receipt; at another,
PLANTATIONOCENE. A third guest asks what keeps the yeast alive, and where
the water comes from.

The Purifier gathers the discarded papers.

“One account is easier to manage,” he says, pushing them beneath the carpet.
\end{ato}

Addressing the bill to humanity does not explain who has ordered, who has
worked, or who is made to pay. Malm and Hornborg challenge accounts that
turn the species into an undifferentiated historical agent.\fonte{malm-hornborg2014}
The names Capitalocene and Plantationocene direct attention towards
historically organised relations of accumulation, labour, land, and
displacement. Haraway's Chthulucene asks how to think and live through
multispecies relations that the figure of a self-sufficient anthropos
fails to describe.\footnote{\refobra{haraway2015}; \refobra{haraway2016}.}
These interventions do not all perform the same task as a stratigraphic
proposal. They also change what is being asked of a name. The problem
for this essay is how to keep those differences active when speaking
about a crisis whose effects and responsibilities are distributed unequally.

\begin{orientacao}{Pós-ato IV --- interrogar o anthropos da conta}
\textbf{Função.} Passar da cronologia às atribuições de agência,
responsabilidade, benefício e exposição.
\begin{enumerate}
\item Desenvolva a crítica ao sujeito humano indiferenciado. Distinga
contribuição causal, benefício econômico, poder de decisão e
vulnerabilidade: essas distribuições não coincidem automaticamente.
Uma responsabilidade planetária compartilhada não implica parcelas iguais.
\item Contextualize \textit{Capitalocene} em relação às contribuições
distintas de Andreas Malm e Jason W. Moore, sem atribuir sua invenção
exclusivamente a Haraway. Para desenvolver Moore, consulte e acrescente
à bibliografia a obra escolhida. Situe \textit{Plantationocene} na
conversa coletiva de Aarhus de 2014 relatada por Haraway em 2015.
\item Apresente o Chthuluceno como exercício de pensamento e composição
multiespécie. Evite tratá-lo simplesmente como outra época a inscrever
na tabela estratigráfica. A pergunta sobre levedura e água desloca
as condições de existência do próprio bar.
\item Reconheça diferenças dentro da literatura do Antropoceno:
Crutzen já registra desigualdades em 2002, e Steffen e colaboradores
desagregam tendências em 2015. A crítica deve especificar os usos
homogeneizadores do nome a que se dirige.
\end{enumerate}
\textbf{Passagem.} Pergunte como a conta única pode parecer evidente:
quais procedimentos deixam relações decisivas fora dos registros?
\end{orientacao}

\begin{ato}{The house rules}
Now the writer watches the Purifier work.

Two signs stand on the counter: NATURE and SOCIETY. The glass goes beneath
the first; the bill beneath the second.

Someone draws a line from the water to the pipe, from the pipe to the
river, from the river to a contract. The Purifier wipes the line away.

“A place for everything,” he says, tapping the house rules.

Behind him, a pipe leaks onto the ledger. Ink runs across the columns.

He fetches a larger cloth.
\end{ato}

Latour's modern Constitution distinguishes Nature and Society while
depending on practices that continuously connect them. Purification and
the proliferation of hybrids belong to the same arrangement.\fonte[pp.~16--17, 35--40]{latour1994}
Read through this argument, the Purifier is a role sustained by the
bar's rules, rather than the personal source of every disturbance.
The pipe, the river, and the contract already act together while the
ledger assigns their effects to separate columns. My working hypothesis
is that contemporary crises are intensified by arrangements that multiply
material mediations while restricting public recognition of the
dependencies and obligations they produce. This is the sense in which
I propose to investigate a failure of communication.

\begin{orientacao}{Pós-ato V --- explicitar a Constituição moderna}
\textbf{Função.} Dar precisão ao mecanismo dramatizado pela limpeza.
\begin{enumerate}
\item Explique conjuntamente tradução/mediação e purificação.
O trabalho moderno produz associações e, ao mesmo tempo, sustenta
a representação de domínios separados. Os híbridos não começam a
existir quando aparecem na contabilidade; sua eficácia precede o
reconhecimento que lhes é concedido.
\item Descreva a alternância entre transcendência e imanência nas
garantias da Constituição. Não use \textit{imanência} como sinônimo
de inércia ou passividade: o problema aqui é a possibilidade de
tratar dependências como exteriores às decisões que as mobilizam.
\item Faça do Purificador uma função institucional reproduzível,
inclusive pelo escritor. Sustente a hipótese comunicacional com
processos materiais: classificação de danos, contratos, infraestruturas,
instrumentos e modalidades de representação.
\item Aproxime essa análise do deslocamento, em \textit{Onde Aterrar?},
de um cenário supostamente estável para um Terrestre que participa
da ação. Localize a discussão no capítulo 9 da edição utilizada:
\refobra{latour2020}.
\end{enumerate}
\textbf{Reserva para a parte II.} A analogia entre Science-bar e o
Deus barrado exigirá explicação própria. Na tradução de 1994 fornecida,
o subtítulo é \enquote{A quarta garantia: a do Deus suprimido},
pp.~38--39. Distinguir sua elaboração do termo usado pelo tradutor
evita apresentar Science-bar como conceito já formulado por Latour.
\end{orientacao}

\begin{ato}{The room gets into the sentence}
Those nearest the carpet have been coughing for some time.

“Speak clearly,” calls a guest from the raised table.

The Purifier sweeps faster. Beneath the fabric, damp receipts have stuck
to soot, splinters and broken glass. Each stroke sends something through
the weave.

The writer raises his voice, as he did when he arrived.

“How the hell can—”

He coughs. Someone brings water; someone opens a window. The draught shifts
the dust towards another table.

The Purifier writes BACKGROUND NOISE in the minutes.
\end{ato}

The disturbance now interrupts the conditions under which a sentence can
be completed. Yet the room was already involved in speech: the elevated
table, the position of the carpet, and the distribution of clean air have
never been neutral. Calling the cough background noise is another
decision about what can count as a contribution. Latour's account of
articulation invites attention to the differences that bodies become
able to register through their encounters.\fonte[pp.~205--13]{latour2004}
Exposure alone, however, guarantees no such learning. A cough may silence
a speaker, and opening a window may transfer an injury. The task is to
follow how damaged conditions of speech are produced, sustained, and
potentially transformed.

\begin{orientacao}{Pós-ato VI --- materializar a falha de comunicação}
\textbf{Função.} Mostrar como a purificação agrava as condições corporais
e institucionais da própria descrição.
\begin{enumerate}
\item Reconstrua um circuito: dependências são omitidas; decisões
prosseguem sem responder por seus efeitos; os efeitos se acumulam;
as condições de participação se deterioram; os sinais dessa
deterioração são novamente classificados como ruído. Sustente esse
circuito com um caso delimitado se o argumento assumir alcance empírico.
\item Defina \textit{comunicação} de modo que inclua circulação de
documentos, instrumentos e obrigações, além das palavras. A crise
não será resolvida apenas pela escolha de um vocabulário melhor.
Especifique o que poderia interromper ou alterar o circuito descrito.
\item Retome o \textit{body talk} sem romantizar sofrimento. Aprender
a ser afetado requer dispositivos, relações e possibilidades de
resposta. A tosse não é, por si só, uma revelação epistemológica.
\item Destaque a desigualdade espacial: alguns já tossiam antes do
escritor. O gesto de abrir a janela mostra que uma resposta local
pode redistribuir o dano. A poeira é heterogênea; não substitui todos
os não humanos por uma única matéria indiferenciada.
\end{enumerate}
\textbf{Passagem.} O escritor precisa alterar sua relação com a pergunta
e com o nome que aprendeu a usar. Aumentar o volume não basta.
\end{orientacao}

\begin{ato}{A word held at a distance}
The writer takes a napkin and writes Anthropocene.

He looks at the bottle, the sediment, the bill. Then he places quotation
marks around the word.

“Whose word?” asks the bartender.

“One I have learned to use,” he says. “It names the time I live in.
It does not tell you where I am speaking from.”

His hand leaves a grey print beside the sentence.

The bartender turns the napkin towards the others. There is still room
to write.

Under the table, the broom keeps moving.
\end{ato}

I retain the word \anthropocene, but I can no longer treat it as a name
that speaks from nowhere. Its quotation marks register a distance
between inhabiting the transformations it gathers and consenting to
the position from which they are named. That distance offers no
exemption: the hand drawing the marks is implicated in the room it
describes. Drawing on Manuela Carneiro da Cunha's distinction between
culture and culture in quotation marks, I propose to investigate this
distance as a problem of situated borrowing and translation.\fonte[chap.~18]{cunha2017}
How might we use an inherited name while making perceptible the
mediations through which it travels, the differences it gathers, and
those it leaves unheard?

\begin{orientacao}{Pós-ato VII --- tornar necessária a segunda parte}
\textbf{Função.} Introduzir o problema das aspas e delimitar o que ainda
precisa ser desenvolvido em \textit{Quotation marks and Science-bar}.
\begin{enumerate}
\item Explique o estrangeirismo como relação de enunciação: a palavra
é recebida de uma posição que quem a emprega não assume integralmente.
Isso pode ocorrer dentro da mesma língua. Identifique a posição
de quem escreve, sem reivindicar uma voz indígena, universal ou
exterior à modernidade que o texto não possa sustentar.
\item Articule essa operação com o capítulo 18 de Cunha,
\enquote{\enquote{Cultura} e cultura: conhecimentos tradicionais e
direitos intelectuais}. A seção sobre conversas de \enquote{cultura},
\textit{kastom} e termos de empréstimo é especialmente pertinente:
o empréstimo envolve tanto sentidos quanto indicações sobre como
interpretar a fala. Explique onde sua proposta prolonga essa análise.
\item Diferencie essa função das aspas de simples ironia ou dúvida
sobre a realidade física das crises. O termo já aparece entre aspas
em documentos científicos: a contribuição do ensaio está na função
que lhes atribui, não na invenção do sinal tipográfico.
\item Mantenha a marca cinza da mão e o movimento da vassoura.
As aspas não suspendem a crise nem absolvem quem escreve; abrem
uma exigência de tradução e transformação das condições de fala.
\end{enumerate}
\textbf{Abertura.} A parte II deverá explicar como Science-bar e
Cecinepá respondem a essa exigência. Reserve a elaboração de futuro
ancestral, tempo espiralar e mundo comum para a parte III, preservando
as diferenças entre seus interlocutores.
\end{orientacao}

\section{Quotation marks and Science-bar}
\label{sec:quotation}

\begin{orientacao}{Desenvolvimento posterior --- sete lugares reservados}
Esta parte recebe a pergunta deixada pelo guardanapo. Os blocos abaixo
são espaços de planejamento, não novos atos já redigidos. A estrutura
preserva sua proposta de sete atos por parte. Substitua cada bloco por
um ambiente \texttt{ato} e escreva o trecho acadêmico correspondente.
\end{orientacao}

\atoplanejado{I}{Retomar a circulação do guardanapo. Desenvolver o
estrangeirismo enunciativo e a mudança de registro em diálogo com
\refobra{cunha2017}. Especificar quem toma emprestado o nome e quais
obrigações a tradução produz.}

\atoplanejado{II}{Alterar uma regra do bar para interrogar a pretensão
de uma Ciência que encerra a discussão. Apresentar Science-bar como
proposta do ensaio e precisar sua relação com a quarta garantia da
Constituição moderna: \refobra{latour1994}, pp.~38--39. Explicar
também os limites dessa analogia.}

\atoplanejado{III}{Dar corpo e localização a Cecinepá. Explicitar se é
um lugar existente, uma experiência, um dispositivo de escrita ou
uma combinação desses modos. Desenvolver sua condição de heterotopia
com \refobra{foucault2013}, sem fazer de todo espaço alternativo uma
heterotopia automaticamente.}

\atoplanejado{IV}{Fazer aparecer uma situação em que coisas, pessoas,
palavras e obrigações não possam ser inventariadas isoladamente.
Definir a proposta de \textit{total ontological fact} e distinguir
seu alcance do fato social total de Mauss. Usar o \textit{Ensaio
sobre a Dádiva}, em \refobra{mauss2017}, como interlocução explícita.}

\atoplanejado{V}{Experimentar o arcabouço semiótico com um caso.
Mostrar como ele descreve mediações e diferenças, quais categorias
emprega e o que permite perceber que a classificação anterior
deixava fora. Tornar o funcionamento de Cecinepá verificável no texto.}

\atoplanejado{VI}{Introduzir uma discordância ou um fracasso do
próprio experimento. Definir como o arcabouço recebe recalcitrância,
assimetria e incomensurabilidade sem pressupor tradução completa
ou consenso de antemão.}

\atoplanejado{VII}{Transformar as condições de participação no bar.
Abrir a passagem para a composição de um mundo comum, sem apresentar
a proposta semiótica como solução suficiente para todas as crises.}

% Seu texto continua aqui. Os titulos dos atos sao opcionais e editaveis.
% \begin{ato}{Titulo do primeiro ato}
% Narrative in English.
% \end{ato}
% Academic prose in English.\fonte[chap.~18]{cunha2017}

\section{A Big Bang towards ancestral futures and the common world}
\label{sec:futures}

\begin{orientacao}{Desenvolvimento posterior --- composição e temporalidade}
O \textit{Big Bang} do título precisa ser definido como operação do
ensaio, caso não se trate do conceito cosmológico. Evite sugerir uma
criação a partir do nada que apagaria histórias, relações e
ancestralidades. Os sete espaços seguintes permanecem editáveis.
\end{orientacao}

\atoplanejado{I}{Retomar o bar transformado. Mostrar que compor o comum
requer trabalho continuado, a partir de dependências já existentes;
definir o sentido do \textit{Big Bang} neste percurso.}

\atoplanejado{II}{Introduzir \refobra{krenak2022}. Desenvolver a
ancestralidade como relação com possibilidades de existência e de
futuro, com atenção à situação e à voz de Krenak. Evitar uma noção
genérica de sabedoria ancestral.}

\atoplanejado{III}{Introduzir \refobra{martins2021}. Relacionar tempo
espiralar, oralitura e corpo-tela às formas do ensaio e à circulação
de sua estória; reconhecer as matrizes afro-diaspóricas específicas
dessa elaboração.}

\atoplanejado{IV}{Fazer os corpos retomarem gestos e palavras do início
sob novas condições. Desenvolver o encontro entre Latour e Martins
sem reduzir uma teoria do corpo à outra. Mostrar convergências,
diferenças e transformações produzidas pela leitura conjunta.}

\atoplanejado{V}{Voltar à conta, ao rio e às condições de respirar.
Especificar obrigações, formas de participação e meios de sustentar
a composição do comum. A proposta precisa alcançar práticas além
da mudança de vocabulário.}

\atoplanejado{VI}{Confrontar um limite da composição: conflito,
perda, assimetria ou entidade que não caiba nos acordos. Distinguir
mundo comum de homogeneização ontológica e de consenso obrigatório.}

\atoplanejado{VII}{Retomar a pergunta inicial e transformar seu modo
de enunciação. Concluir sem apagar a poeira ou encerrar a
responsabilidade pelas mediações. Indicar o que o ensaio tornou
possível e o que permanece como tarefa.}

% Inserir aqui a narrativa e a argumentacao da terceira parte.


\clearpage
\section*{Works Cited}
\begin{orientacao}{Bibliografia de trabalho}
No modo orientado, aparecem as referências preparadas para as três partes.
No modo sem orientações, aparecem somente as obras chamadas por
\texttt{\string\fonte} ou \texttt{\string\refobra} no texto visível.
Confira os localizadores de cada passagem na edição efetivamente utilizada.
Nos livros digitais sem paginação estável, use capítulo ou seção: o número
da página do visualizador de PDF não equivale necessariamente à página da
edição. As entradas estão em ordem alfabética e as notas usam autor e
título abreviado, conforme os exemplos da revista.
\end{orientacao}
\imprimirreferencias
\end{document}
